problem:  
Let \( (M^n,g) \) be a closed, simply connected Riemannian manifold with sectional curvature \( K_g \ge 1 \). Then automatically \( \mathrm{Ric}_g \ge (n-1)g \), and hence \( \lambda_1(M,g) \ge n \) by Lichnerowicz–Obata, with equality if and only if \( (M,g) \) is isometric to the unit round sphere \( S^n \).  

**Question (Quantitative spectral–geometric stability for the sphere):**  
Does there exist, for each \( n \), a constant \( C_n>0 \) such that  
\[
|\operatorname{Vol}(S^n) - \operatorname{Vol}(M,g)| \le C_n \,|\lambda_1(M,g) - n|?
\]
for every closed, simply connected manifold \( (M^n,g) \) with \( K_g \ge 1 \)?  
Equivalently, does *spectral near-equality* to the sphere force *quantitative volume closeness* at a **linear rate** in the spectral defect?  

More geometrically: if \( K_g \ge 1 \) and \( \lambda_1(M,g) \le n+\varepsilon \), can one guarantee  
\[
|\operatorname{Vol}(S^n) - \operatorname{Vol}(M,g)| \le C_n \varepsilon
\]
and Gromov–Hausdorff closeness \( d_{GH}(M,S^n) \le f_n(\varepsilon) \) for a universal function \( f_n \) vanishing linearly as \( \varepsilon\to0 \)?  

Why is it a "good" problem:  
This refined statement isolates a single, mathematically coherent phenomenon—the *quantitative stability* of spectral and geometric sphere rigidity under sectional curvature control. It represents a natural and profound extension of classical theorems:

- **Lichnerowicz–Obata** gives a sharp qualitative rigidity result,
- **Toponogov** and **Myers–Bonnet** control geometry by curvature,
- yet no known technique quantifies how *close* a positively curved manifold must be to the sphere once \( \lambda_1 \) or volume nears the spherical value.

The conjectured inequality provides a **precise bridge** between *spectral analysis* and *metric geometry*: it asks whether nearness of analytic data (the Laplace spectrum) enforces nearness of global shape and measure. Establishing or disproving such a linear relation would reveal the analytically measurable stiffness of positively curved spaces, paralleling Cheeger–Colding’s results for lower Ricci bounds but at the more delicate sectional curvature level.  

The statement is elegant—just a relationship between volume and the first eigenvalue—but resolving it would require new techniques connecting curvature comparison, geometric measure estimates, and spectral perturbation theory. Such a problem is both fundamental and far-reaching, qualifying as a genuinely "good" research-level problem.